\section{Introduction}
\label{sec:introduction}


While Large Language Models (LLMs) have been widely accepted for performing various reasoning tasks~\citep{kojima2022large,wang2022self,lightman2023let,wei2022chain}, recent studies have revealed that they still struggle in multi-turn and long-horizon settings~\citep{laban2025llms}: models tend to propagate previous errors and incorrect assumptions rather than revise them. A key limitation of their autoregressive (AR) paradigm is that local errors require regenerating the entire sequence, incurring unnecessary computation. Moreover, incorrect intermediate content persists in the context, occupying valuable capacity and contaminating subsequent reasoning.

In contrast, Mask Diffusion Models (MDMs)~\citep{austin2021structured, sahoo2024simple, nie2025large, xin2025lumina} provide a fundamentally different generation mechanism that naturally supports localized revision. Through iterative masked updates, MDMs can keep the current context fixed while resampling uncertain tokens, avoiding the need to regenerate the entire sequence. This native remasking mechanism reveals a potential advantage of MDMs for reasoning: if local errors can be revised in place, the model could maintain a cleaner intermediate context, reduce error contamination, and save computation. In this sense, MDMs offer a possible path toward localized self-correction that is difficult to realize in AR generation. 

However, realizing this potential requires moving beyond the passive remasking behavior of existing MDM decoding. Standard MDM generation still follows an absorbing Markovian denoising process: once a token is confidently denoised, it is fixed. Existing remasking strategies focus on re-masking low-confidence tokens, but the model is unable to actively revisit and correct previously committed predictions. We argue that for MDMs to support reasoning-level self-correction and provide a distinct form of test-time scaling from AR models, the model must learn to identify unreliable predictions and actively revise them during generation. We propose \textbf{\methodfull}, where masking is formulated as an internal decision process driven by the model's uncertainty, enabling it to selectively refine already denoised tokens during generation.

In this work, we activate \textbf{\methodfull as a first-class capability} of MDMs by making it a central focus of post-training. We introduce a lightweight training paradigm that enables \textbf{self-initiated, context-aware revision} without architectural modifications, together with an effective data generation strategy that produces stable training signals aligned with the model’s native output distribution. Our results show that this latent capability can be reliably unlocked through appropriate training. More broadly, reflective masking establishes a form of test-time scaling unique to MDMs, where additional computation is allocated to selective revision rather than forward expansion. We argue that reasoning in MDMs should be viewed not as forward generation, but as iterative state refinement through explicit self-correction and revision.

With multi-turn reflective masking, a key capability still missing relative to autoregressive reasoning is the ability to extract insights from the growing context of historical generations. To bridge this gap, we introduce \textbf{History Reference}, a parameter-free mechanism that enables MDMs to maintain a stateful view of their denoising trajectory by preserving intermediate decoding states. Unlike AR models, which encode history implicitly in the input context, History Reference allows MDMs to explicitly leverage past predictions. This mechanism introduces a temporal dimension orthogonal to the current text context, guiding reflective updates and improving consistency while reducing repeated errors during iterative refinement. Importantly, History Reference requires no additional learnable parameters or external memory, making it efficient and readily applicable to existing MDMs.

We evaluate our approach across a spectrum of generation tasks with varying levels of guidance, demonstrating consistent improvements and strong generality. We begin with image editing, where rich instructions specify both where and how to modify the input. We then consider Sudoku, a structured reasoning task that requires the model to identify and correct erroneous entries. Finally, we study text generation tasks, where supervision is minimal and no direct hints about the final answer are provided. Across all three settings, \textbf{\methodfull} consistently improves performance. For tasks that require autonomous exploration, \textbf{History Reference} proves particularly effective in guiding the model’s reasoning and iterative revision. Together, these results suggest that reflective masking provides a general mechanism for enabling reasoning through explicit revision in mask-based generation. \looseness-1

Our contributions can be summarized as follows:
\begin{itemize}[leftmargin=*]
    \item We identify \textbf{\methodfull} as a latent capability of mask diffusion models, and propose a new perspective that frames generation as an iterative process of self-initiated revision, introducing a new dimension of reasoning based on explicit modification of prior outputs.

    \item We propose a lightweight training paradigm, associated with an effective and scalable data pipeline, to activate \methodfull as an intrinsic skill, enabling context-aware and adaptive revision without architectural changes. We demonstrate its effectiveness across diverse downstream tasks and modalities, including text generation, Sudoku, and image editing, highlighting strong generality.

    \item We introduce \textbf{History Reference}, a parameter-free mechanism that enables MDMs to maintain a stateful view of their denoising trajectory by incorporating intermediate decoding history. It improves revision consistency and avoids repeating past errors during iterative refinement.\looseness-1
\end{itemize}